\documentclass[11pt]{article}
\usepackage[margin=1in]{geometry}
\usepackage{amsmath, amssymb, amsthm, mathtools}
\usepackage{stmaryrd}
\usepackage{hyperref}
\usepackage{listings}
\usepackage{xcolor}
\usepackage{enumitem}
\usepackage{titlesec}
\usepackage{fancyhdr}
\usepackage{quoting}

\definecolor{rhonavy}{HTML}{1A1A2E}
\definecolor{rhosky}{HTML}{3FA9F5}
\definecolor{rhoyellow}{HTML}{F3D630}
\definecolor{rhosage}{HTML}{8BB999}

\hypersetup{
  colorlinks=true,
  linkcolor=rhonavy,
  urlcolor=rhosky,
  citecolor=rhonavy
}

\titleformat{\section}{\Large\bfseries\color{rhonavy}}{\thesection}{1em}{}
\titleformat{\subsection}{\large\bfseries\color{rhonavy}}{\thesubsection}{1em}{}
\titleformat{\subsubsection}{\normalsize\bfseries\color{rhonavy}}{\thesubsubsection}{1em}{}

\pagestyle{fancy}
\fancyhf{}
\fancyhead[L]{\small\color{rhonavy}\textsc{Position Paper}}
\fancyhead[R]{\small\color{rhonavy}\textsc{Typed Value}}
\fancyfoot[C]{\thepage}
\renewcommand{\headrulewidth}{0.4pt}
\renewcommand{\headrule}{\hbox to\headwidth{\color{rhosky}\leaders\hrule height \headrulewidth\hfill}}

\lstdefinelanguage{rholang}{
  morekeywords={new, in, for, contract, match, if, else, true, false, Nil, send, receive, let, return, type, def},
  sensitive=true,
  morecomment=[l]{//},
  morestring=[b]"
}
\lstset{
  basicstyle=\ttfamily\small,
  keywordstyle=\color{rhonavy}\bfseries,
  commentstyle=\color{gray}\itshape,
  stringstyle=\color{rhosky},
  numbers=left,
  numberstyle=\tiny\color{gray},
  frame=single,
  rulecolor=\color{rhosky!50},
  breaklines=true,
  showstringspaces=false,
  columns=fullflexible,
  framexleftmargin=4pt,
  framexrightmargin=4pt,
  framextopmargin=4pt,
  framexbottommargin=4pt
}

\newenvironment{position}
  {\begin{quoting}[indentfirst=true, leftmargin=0pt, rightmargin=0pt, vskip=0.5em]
   \noindent\color{rhonavy}\rule{4pt}{1em}\hspace{0.5em}\bfseries\itshape}
  {\end{quoting}}

\newcommand{\position}[1]{%
  \medskip
  \begin{quoting}[indentfirst=false, leftmargin=1em, rightmargin=1em]
  \noindent\textcolor{rhosky}{\rule{3pt}{1.2em}}\hspace{0.5em}\textbf{\textit{#1}}
  \end{quoting}
  \medskip
}

\newcommand{\Sig}{\mathsf{Sig}}
\newcommand{\Pay}{\mathsf{Pay}}
\newcommand{\Value}{\mathsf{Value}}
\newcommand{\Contract}{\mathsf{Contract}}

\title{
{\Large \textsc{Position Paper}}\\[0.5em]
{\Huge\color{rhonavy} Typed Value}\\[0.4em]
{\Large Currency as Proof-Bearing Process}\\[0.6em]
\large Arbitrage as Confession of a Gameable Score;\\
\large the Reduction to Practice in MeTTaIL
}
\author{L.G.~Meredith\\\small F1R3FLY Industries\\[0.25em]\small \textit{in dialogue with Claude (Anthropic)}}
\date{\today}

\begin{document}
\maketitle

\begin{abstract}
\noindent We argue that fiat currency, like the typeless-scalar reputation systems prevalent on existing platforms, is a gameable scoring system, and that the persistent multi-trillion-dollar arbitrage industry is the deductive evidence of its gameability rather than incidental friction in an otherwise sound measurement. A scalar that does not exhibit the process it witnesses cannot be adjudicated. We extend the typed-judgment framework developed for reputation to value transfer, exhibit a constructive replacement in which payments carry spatial-behavioral types describing the process they witness, and show that this is not speculative: the substrate already runs in MeTTaIL via Simon Peyton Jones-style financial contract combinators, phlogiston-metered valuable friction, and the F1R3FLY shard infrastructure. The position is that typed value is the natural object replacing the fiat scalar, and that the move from platform to polity follows when the unit of account is replaced by an adjudicable judgment.
\end{abstract}

\vspace{1em}

\section{Thesis}

\position{A scalar score that does not exhibit the process it witnesses is not a measurement. It is a commitment to forget. Fiat currency is such a scalar; arbitrage is the proof; typed value is the replacement; and the computational substrate to host typed value already runs.}

The argument develops in three escalating moves.

The first move was established elsewhere: typeless-scalar reputation systems are structurally gameable because the aggregation process the scalar is supposed to witness is unrecoverable from the scalar itself. Disputes cannot be adjudicated, because what is being measured is not exhibited. The cure is to expose the structured object underneath --- signed proofs of capability and attested traces of observation --- and let communities project to scalars only at specific decision points, with the projection itself adjudicable.

The second move is the position of this paper: \emph{fiat currency suffers the same defect, more severely}. More severely because fungibility is the explicit design goal of currency in a way it is not for reputation; because the stakes are higher; because the manipulation infrastructure is more developed; and because the gap between the scalar and the underlying process is harvested as a primary industry rather than treated as a defect. Arbitrage is to currency what brigading is to platform reviews: not a regrettable side-effect of an otherwise good measurement, but the deductive symptom that no coherent measurement is taking place.

The third move is that this is reducible to practice. We do not need a hypothetical computational economy to host typed value. We need exactly the substrate that has been engineered at F1R3FLY: rho-calculus communication metered by phlogiston, MeTTaIL as the language-definition layer, and Simon Peyton Jones-style financial contract combinators already implemented in rholang and extended to loans, Dutch auctions, and power-market instruments. The position paper's argument is therefore not ``we should build this'' but ``the construction exists; here is why it is the correct response to the structural defect of fiat.''

\section{The reputation precedent}

\subsection{What was wrong with the scalar}

In the typed-judgment reputation framework, an agent's capability is adjudicated by two coupled processes: a distributed proof system over signed judgments
\[
\Gamma \vdash \Sig(A) : T
\]
and a communal labelled transition system assembled from attested communication events. Reputation is the structured pair of provability and refutation, projected to scalars only by community-specific aggregation functions whose choice is itself open to adjudication.

The structural defect of the alternative --- the typeless scalar familiar from review platforms, credit scores, and karma systems --- is that the aggregation function is not exhibited at the rating site. The output is a number; the process that produced the number is opaque; and the dispute resolution is therefore impossible. A reviewer cannot challenge a 4.3-star rating because what 4.3 stars \emph{is}, in terms of the underlying behavior it claims to summarize, has been discarded.

\position{The diagnosis: a scalar without a process is not measurement; it is the commitment to forget what one was measuring.}

\subsection{Why the cure generalizes}

The cure for the reputation case was not the elimination of summary judgments but the \emph{recovery of process}. Replace the scalar with a structured judgment about whether an agent inhabits a typed pattern of behavior; let the community accumulate signed proofs and attested observations against the judgment; expose the entire object; project to scalars at decision points only, with the projection itself open to scrutiny.

The structural ingredients of this cure --- cryptographic identity, type-indexed claims, signed evidence with provenance, evidence-indexed modal grading --- are not specific to reputation. They are the general apparatus for adjudicating claims about opaque agents. Wherever a scalar is currently doing the work of summarizing such claims, the same cure applies. \emph{Currency is the next case.}

\section{Currency as a scoring system}

\subsection{The analogy made precise}

A reputation score and a fiat unit of account perform structurally identical functions:
\begin{itemize}[leftmargin=*]
\item Both are scalars produced by an aggregation process that is not exhibited at the point of use.
\item Both circulate as if they denoted a coherent underlying quantity.
\item Both are governed by an authority (platform / central bank) whose aggregation function is not open to adjudication by holders of the score.
\item Both are subject to manipulation by actors who understand the gap between the scalar and the underlying process better than the typical holder does.
\end{itemize}

The dollar held in a wallet is the karma of an account on a forum. Both are tokens of standing in a system, denominated in units whose meaning is presupposed rather than exhibited.

\subsection{Why currency is the more severe case}

The defects of the reputation scalar are bounded in three ways the fiat scalar's are not.

First, \emph{stakes}. A miscalibrated reputation score affects access to platform features; a miscalibrated currency unit allocates the resources of an entire economy. The cost of gameability scales with the social weight of the system.

Second, \emph{fungibility as design goal}. Reputation scores at least nominally aspire to track the agent and the context (Alice on Topic X is differentiated from Alice on Topic Y in well-designed systems). Currency aspires to maximum fungibility: a dollar is a dollar regardless of where it came from, what it was paid for, or what it will buy. This is not a defect of implementation but the explicit ambition of the design. The currency unit is constructed to forget.

Third, \emph{the manipulation surface is mature and lucrative}. Review-brigading on a platform might generate cents per gamed rating. Currency manipulation generates basis points on trillions of dollars of notional value, supports an entire industry, and is sufficiently lucrative to drive the construction of dedicated infrastructure --- co-located trading systems, high-frequency networks, regulatory-arbitrage corporate structures. The gap between the scalar and the process is so reliably present and so reliably profitable that harvesting it is an asset class.

\position{Fiat currency is not a misimplemented good measurement. It is a precisely-implemented forgetful one. The forgetting is the design.}

\section{Arbitrage as confession}

\subsection{The deductive argument}

If currency were a coherent measurement of value --- if a dollar denominated the same quantity regardless of where, when, and against what it was transacted --- then the following would hold by definition:

\begin{quote}
For any commodity, service, or asset $X$, the price of $X$ in dollars at any market $M$ would be identical to the price of $X$ in dollars at any market $M'$, modulo costs of transfer that themselves denominate the same quantity.
\end{quote}

This is not a hypothesis. It is a consequence of what it means for a measurement to be coherent. Two measurements of the same quantity, by the same instrument, must agree.

\emph{Arbitrage exists}: profitable opportunities exist to buy $X$ at price $p_1$ on market $M_1$ and sell $X$ at price $p_2 > p_1 + \text{transfer}$ on market $M_2$. These opportunities are not occasional; they are persistent; they support a multi-trillion-dollar industry; and their elimination at any single moment is undertaken by the same arbitrageurs whose continued existence proves the elimination is incomplete.

The conjunction of the consequence and the empirical fact is the deductive argument: \emph{the dollar does not denominate the same quantity in $M_1$ and $M_2$}. What looks like a unit of account is locally inconsistent. The scalar fails to track a single coherent underlying process.

\position{If currency were a coherent measurement, arbitrage would be impossible. The persistent presence of arbitrage opportunities is logically equivalent to the local inconsistency of the scalar.}

\subsection{Variants of the same confession}

The taxonomy of arbitrage strategies is the taxonomy of the ways the scalar fails to track the process.

\paragraph{Spatial arbitrage.} The same asset priced differently in geographically separated markets. The scalar's unit is supposed to be jurisdiction-invariant; it is not.

\paragraph{Temporal arbitrage.} Carry trades, basis trades, forward arbitrage. The scalar at $t$ does not coherently denominate against the scalar at $t'$; the difference is a profit opportunity.

\paragraph{Regulatory arbitrage.} Same transaction, different jurisdictions, different tax/regulatory treatment, different effective price. The scalar erases the regulatory context, but the context is doing work the scalar pretends to capture.

\paragraph{Latency arbitrage.} High-frequency exploitation of microsecond-scale price differences across venues. The information state of the scalar is not synchronized with the information state of the underlying process.

\paragraph{Statistical arbitrage.} Systematic profit extraction from pricing inefficiencies whose existence presupposes that prices are not informationally complete summaries of underlying processes.

\paragraph{Foreign exchange arbitrage.} The classical triangular case: three currencies whose pairwise exchange rates do not compose to unity. The scalar fails even at the level of internal arithmetic consistency.

\paragraph{Monetary policy arbitrage.} Anticipation of central bank action allows positions sized to capture the redistribution of purchasing power that follows. The scalar's value is altered by an exogenous process whose timing is partially predictable; predictability is profit.

Each is a window onto the same underlying fact. The dollar does not measure a single thing. Where the gap between any two of its instances is exploitable, the scalar has confessed that no coherent quantity is being denominated.

\subsection{The forgetting that produces the gap}

Why is the gap so reliably present? Because the construction of the scalar is the systematic erasure of the information that would close the gap. The dollar paid for a derivative trade, the dollar paid for groceries, the dollar paid for a healthcare procedure, the dollar paid as a campaign contribution --- these are denominated identically and circulate identically. The processes they witness are radically different; the unit of account is the commitment to treat them as the same.

\position{The gap that arbitrage exploits is not a defect introduced by friction in an otherwise coherent system. It is the necessary consequence of a system whose design principle is the erasure of process.}

This is the deeper move beneath the arbitrage argument. The arbitrageur is not merely a thief from a broken till. The arbitrageur is the consequence of the till's design: a system that erases process \emph{must} produce exploitable gaps, because process is what would otherwise close them.

\section{What money systematically erases}

To make the framework move precise, it is useful to enumerate what the fiat scalar erases that a typed alternative would preserve.

\paragraph{Provenance.} A dollar carries no record of its origin. A dollar created by central bank asset purchases, a dollar earned by labor, a dollar laundered through three shell companies are scalar-identical. The community has no structural way to query a unit's history.

\paragraph{Process.} A dollar carries no record of what it is a payment for. The receiving party may know; the system does not. ``GDP grew by 3\%'' is the aggregate scalar in the same units regardless of whether the growth was driven by hurricane reconstruction or pharmaceutical research.

\paragraph{Jurisdiction.} The dollar is nominally domestic but circulates globally. Its effective regulatory context at any moment is a function of where it is held, which the scalar does not encode.

\paragraph{Counterparty risk.} The dollar in a deposit account and the dollar in a mattress are scalar-identical, but they represent different claims against different counterparties with different solvency properties. The aggregation collapses the distinction.

\paragraph{Capability claims.} A dollar promised in a contract and a dollar already in hand are scalar-identical for purposes of bookkeeping, but represent radically different epistemic states about the world.

\paragraph{Time of creation.} A dollar created in 1971 and a dollar created in 2023 are scalar-identical at any present moment, but they were issued against different baskets of underlying assets, by different regimes, under different policy commitments.

Each of these axes is information the underlying process \emph{has}, which the scalar \emph{discards}. The discarded information is then exactly what arbitrage, money laundering, regulatory gaming, and monetary policy manipulation reconstruct, in private, for private benefit. The information is not absent from the economy; it is absent from the unit of account, which means it is asymmetrically available.

\section{Typed value: the framework}

\subsection{The proposal}

Replace the scalar with a structured judgment. A payment is no longer a quantity in a fungible unit; it is a typed object of the form
\[
\Gamma \vdash \Sig(A) \xrightarrow{\pi} \Sig(B) : \Pay(\tau)
\]
where $\pi$ is the payment event, $\tau$ is a spatial-behavioral type describing the process the payment witnesses (the goods or services exchanged, the contractual conditions, the regulatory context), and $\Gamma$ records the environmental dependencies on which the payment relies.

The typed payment inherits, directly, the entire apparatus of the reputation framework:
\begin{itemize}[leftmargin=*]
\item Cryptographic identity in the participant positions.
\item Signed claims as the leaves of distributed proofs.
\item Attested traces of communication events building the communal LTS.
\item OSLF-generated Hennessy--Milner logic for adjudication.
\item Evidence-indexed modal grading for refutation.
\item Anti-scalarization: scalar summaries only at specific decision points, with the projection itself adjudicable.
\end{itemize}

The dollar denomination of a payment, if it survives at all, is a projection chosen by a specific community for a specific decision. It is not the unit of account.

\subsection{What typed value resists}

The failure modes of the fiat scalar fall out as resistances of the typed alternative.

\paragraph{Arbitrage resistance.} If $\Pay(\tau)$ on market $M_1$ trades at price $p_1$ and $\Pay(\tau)$ on market $M_2$ trades at price $p_2 \neq p_1$, the type system flags a contradiction: the same typed object cannot have two prices. Resolution requires either (a) explicit refinement showing that the types are actually distinct ($\Pay(\tau_1)$ vs $\Pay(\tau_2)$), making the difference adjudicable, or (b) price convergence. The arbitrageur's traditional opportunity becomes a type-checking audit, with the gap exposed rather than harvested.

\paragraph{Provenance preservation.} The provenance of a typed value is structurally tracked through the proof tree that constructed it. Money laundering does not eliminate provenance --- it would require constructing a derivation whose leaves are forged signatures, which is bounded by the cryptographic assumptions.

\paragraph{Process preservation.} The type $\tau$ is exactly what the scalar erased. A payment for healthcare and a payment for a derivatives trade are typed-distinct objects. The community can query their distribution, audit their aggregation, and resist the kinds of accounting that produce broken-window-fallacy GDP.

\paragraph{Capability claim preservation.} A typed payment that is conditional on capability ($\Sig(B) : T_{\text{service}}$ as a hypothesis in $\Gamma$) carries that conditionality explicitly. It cannot be substituted for an unconditional payment without an explicit type cast.

\paragraph{Counterparty visibility.} The participant identities are first-class in the judgment. The dollar-in-a-deposit-account and dollar-in-a-mattress collapse cannot happen, because the typed objects record their counterparty structure.

\position{Typed value is what the dollar would be if the dollar carried the information the economy already has.}

\section{Reduction to practice}

The position would be vacuous if the substrate did not exist. It does.

\subsection{SPJ combinators as the typed-money primitive}

Simon Peyton Jones, Jean-Marc Eber, and Julian Seward published, in 2000, a small algebra of typed combinators in which the entire landscape of financial contracts --- bonds, swaps, options, exotic derivatives --- can be expressed compositionally. The combinators are typed; their composition is type-checked; the semantics is given as a function from contracts into expected-value processes over time.

The primitive combinators include:
\begin{lstlisting}[language=rholang]
zero      : Contract
one       : Currency -> Contract
give      : Contract -> Contract
and       : Contract -> Contract -> Contract
or        : Contract -> Contract -> Contract
cond      : Obs Bool -> Contract -> Contract -> Contract
scale     : Obs Double -> Contract -> Contract
when      : Obs Bool -> Contract -> Contract
anytime   : Obs Bool -> Contract -> Contract
until     : Obs Bool -> Contract -> Contract
truncate  : Time -> Contract -> Contract
\end{lstlisting}

A zero-coupon bond is \texttt{when (at maturity) (scale (konst face) (one currency))}. A European call is \texttt{when (at strike\_date) (or zero (and (give (scale (konst strike) (one usd))) (one asset)))}. And so on through the entire taxonomy.

\position{The SPJ algebra exhibits, constructively, what a typed money looks like. The fact that this has been available in the literature for a quarter-century without being adopted at the level of the unit of account is a function of substrate, not of theory.}

\subsection{Hosting in rholang/MeTTaIL}

The SPJ combinators have been implemented in rholang at F1R3FLY, with extensions to:
\begin{itemize}[leftmargin=*]
\item Loan types (with explicit capability dependencies on borrower and collateral attestations);
\item Dutch auctions (as time-truncated choice combinators with explicit phlogiston metering);
\item Power-market instruments (with grid-state observables and physical-delivery constraints typed at the combinator level).
\end{itemize}

The combinators inhabit rholang processes; contracts are communicated as terms in the MeTTaIL-defined language; the type system is the spatial-behavioral type system of the rho calculus; and the metering of contract execution is the phlogiston metering of computation. A typed payment is a rho process running on the F1R3Node shard, signed by its participants, with its type checked against the OSLF-generated logic, and its provenance preserved structurally.

\begin{lstlisting}[language=rholang,caption={A typed payment in MeTTaIL, schematic.}]
new payCh, settleCh in {
  contract payCh(@from, @to, @typedContract, @sigFrom, @sigTo) = {
    if (verify(sigFrom, from, typedContract) and
        verify(sigTo, to, typedContract) and
        typecheck(typedContract, contextOf(from))) then {
      executeContract(typedContract, from, to) |
      attestEvent(payment(from, to, typedContract))
    } else Nil
  }
}
\end{lstlisting}

The point is not the toy contract code, which is illustrative. The point is that \emph{nothing in this construction is hypothetical}. The execution semantics is the rho calculus, deployed; the typing is the spatial-behavioral discipline, implemented; the metering is phlogiston, in use; the signature discipline is the standard cryptographic apparatus; and the language definition is MeTTaIL, with the BNFC grammar already written for rholang 1.4.

\subsection{Phlogiston as the valuable-friction primitive}

The standard objection to typed money is that fungibility is what makes currency useful, and that typed currency reintroduces the friction that fiat eliminated.

The response is that \emph{this friction is informational}. The reason a dollar circulates frictionlessly is the same reason it forgets. The friction that typed value reintroduces is the friction of carrying provenance, type, and capability constraints across the transaction. This friction is what closes the arbitrage gap, what resists laundering, what makes monetary policy adjudicable rather than imposed.

The principle is already operative in the F1R3FLY architecture as \emph{valuable friction}. Phlogiston-metered computation is friction that prices the work being done; it is not eliminated but priced and exposed. Typed value extends the same principle from computation to value transfer: friction is not the enemy of utility, friction is the price of adjudicability.

\position{Frictionless money is forgetful money. Valuable friction is the price of remembering. The economy that pays this price gains the capacity to adjudicate its own transactions.}

\section{From economy to polity}

\subsection{Typed agents in a typed market}

When value is typed, the actors who hold and exchange it are no longer abstract market participants whose internal structure is opaque. They are typed agents whose capability claims are subject to adjudication. The market is no longer the institution of anonymous exchange; it is the venue in which typed agents adjudicate typed claims through typed value transfers.

The combination of typed reputation (the earlier framework) and typed value (this paper) yields a substrate in which:
\begin{itemize}[leftmargin=*]
\item Agents make capability claims of the form $\Gamma \vdash \Sig(A) : T_{\text{service}}$.
\item Value transfers are typed payments $\Sig(A) \xrightarrow{\pi} \Sig(B) : \Pay(\tau)$.
\item Contracts compose payments and capabilities into typed protocols.
\item The community adjudicates conformance to all of the above.
\end{itemize}

This is no longer an economy in the standard sense. It is a polity --- a body of typed agents engaged in adjudicable exchange under shared rules of inference, with the rules themselves open to adjudication.

\subsection{The community is the moat}

The strategic implication, developed in the companion ``Community Is the Moat'' position, is that platforms whose value derives from network effects on opaque scalar systems are exposed in a way that platforms whose value derives from adjudicable typed exchange are not. The moat is not the scalar score or the proprietary algorithm; it is the community of typed agents whose adjudication of typed claims is what gives the platform its economic substance.

Web 2.0 platforms accumulate engagement on opaque metrics. Their value is captured by whoever controls the aggregation. Web 3.0, in the form of an adjudicable polity over typed value, captures value at the level of the community whose adjudication is the substance of the system. The shift is from platform-as-scalar-aggregator to polity-as-typed-adjudicator.

The F1R3FLY property portfolio --- F1R3Sky, F1R3Eats, F1R3Tunes, F1R3Docs, F1R3Games --- has been engineered specifically to host this shift. Each property is a domain in which the substrate makes typed exchange feasible; each is a venue in which the unit of account can be replaced by an adjudicable judgment, with the community as the adjudicator.

\section{What follows}

The position paper closes with a programmatic statement and recommendations.

\subsection{Programmatic statement}

\position{Fiat currency is a typeless scalar score whose gameability is exhibited by the persistent multi-trillion-dollar arbitrage industry. The defect is not implementational but constitutional: the unit is constructed to erase process. The replacement is typed value, in which payments are structured judgments under an adjudicable type discipline. The substrate to host typed value exists, in production, on the F1R3Node shard, with SPJ-style financial combinators implemented in rholang, phlogiston-metered valuable friction at the computational layer, and MeTTaIL providing the language-definition primitives. The move from platform to polity is the move from scalar aggregation to communal adjudication of typed claims.}

\subsection{Recommendations}

\begin{enumerate}[leftmargin=*]
\item \textbf{Treat the SPJ combinator suite as currency, not as financial product.} The combinators are not a library for hedge funds; they are the constructive proof that typed value is feasible. Frame them accordingly in F1R3FLY documentation, investor materials, and standards proposals.

\item \textbf{Make arbitrage the inversion test.} For any proposed financial infrastructure --- new exchange, new derivative, new payment rail --- ask: does the proposal eliminate arbitrage by construction, or does it create new arbitrage opportunities? Eliminate-by-construction is the marker of typed value; create-arbitrage is the marker of yet another scalar.

\item \textbf{Position phlogiston as the price of remembering.} The valuable-friction story is the answer to the fungibility objection. Develop the rhetorical and technical articulation of valuable friction as the informational counterpart of computational metering.

\item \textbf{Deploy in the property portfolio.} F1R3Tunes, F1R3Eats, and F1R3Games are the testbeds for typed payments. Each domain has well-typed exchanges (a song stream, a meal, a move in a game) and the substrate to execute them through typed combinators. Make the in-property currencies typed from inception.

\item \textbf{Treat the VA CRADA as the regulated-domain proof of concept.} Healthcare data exchange is the case where typed value transfer with full provenance, capability constraint, and selective-disclosure is not optional but required. The CRADA is the venue in which typed value is demonstrably necessary; treat it as the existence proof.

\item \textbf{Develop the polity literature.} The intellectual frame --- from platform to polity, from scalar aggregation to typed adjudication --- is the differentiator. Make the philosophical articulation match the technical reduction-to-practice.
\end{enumerate}

\bigskip

\hrule

\bigskip

\noindent\textit{This is the second in a sequence of position papers articulating the typed-judgment programme. The first established the framework for typed reputation; this one extends it to typed value. Subsequent papers will address the consensus-synthesis structure of typed-polity governance, the connection to Coalition Logic via Curry--Howard, and the relationship between typed value transfer and the spatial dimension of the rho calculus.}

\end{document}
